% Chapter 4: The binomial model.

\Cref{ch:03-arbitrage} priced claims by static replication: a forward contract, or a
combination of calls and puts, is matched by a portfolio bought today and held to
maturity. A call option cannot be replicated in this way by the stock and the money
market, because its payoff is not an affine function of the stock price. It can be
replicated \emph{dynamically}, by a portfolio that is rebalanced as prices move. The
simplest model in which this can be done is the binomial model of \textcite{crr1979}: in
each period the stock price is multiplied by one of two factors.

Everything in this chapter is finite, and every statement comes down to a finite system
of linear equations or inequalities. Yet the central ideas of the course all appear:
replication, the risk-neutral measure, discounted prices as martingales, the fundamental
theorem of asset pricing, optimal stopping for American options and, with the trinomial
model, what goes wrong when replication fails. \Cref{ch:06-black-scholes-formula} lets the number of periods tend to
infinity and recovers the Black--Scholes formula.

\section{The model}\label{sec:bin-model}

\begin{definition}[The binomial model]\label{def:bin-model}
Fix a number of periods $N\ge1$, an initial stock price $S_0>0$, factors $0<d<u$, an
interest rate $r>-1$ per period, and $p\in(0,1)$.
\begin{enumerate}[label=(\roman*)]
  \item The outcomes are the sequences of $N$ coin tosses, $\Omega=\{H,T\}^N$, written
    $\omega=\omega_1\omega_2\cdots\omega_N$. The tosses are independent, with probability
    $p$ of heads: $\Prob(\{\omega\})=p^{\#H(\omega)}(1-p)^{\#T(\omega)}$, where
    $\#H(\omega)$ and $\#T(\omega)$ count the heads and the tails in $\omega$.
  \item The stock price starts at $S_0$ and is multiplied by $u$ after a head and by $d$
    after a tail: $S_n=uS_{n-1}$ if $\omega_n=H$ and $S_n=dS_{n-1}$ if $\omega_n=T$, for
    $1\le n\le N$.
  \item The money market turns one dollar at time $0$ into $(1+r)^n$ at time $n$; money
    can be lent and borrowed at this rate.
  \item The information at time $n$ is $\F_n=\sigma(\omega_1,\dots,\omega_n)$, generated
    by the first $n$ tosses, with $\F_0=\{\emptyset,\Omega\}$.
\end{enumerate}
The \emph{one-period model} is the case $N=1$.
\end{definition}

The filtration $(\F_n)$ (\cref{def:pr-filtration}) is generated by the partition of
$\Omega$ into the $2^n$ sets on which $\omega_1\cdots\omega_n$ is fixed, so a random
variable is $\F_n$-measurable if and only if it depends on $\omega$ only through the first
$n$ tosses (\cref{def:pr-measurable}); we then write it $Y(\omega_1\cdots\omega_n)$. For
instance $S_n=S_0u^{H_n}d^{n-H_n}$, with $H_n$ the number of heads among the first $n$
tosses. Since $d<u$, the price moves reveal the tosses: $\F_n=\sigma(S_1,\dots,S_n)$.
Every outcome has positive probability, so ``almost surely'' means ``for every $\omega$'',
and conditional expectations given $\F_n$ are defined at every $\omega$ by the formula for
finite partitions (\cref{def:pr-condexp-discrete}).

A \emph{claim} (or derivative) with maturity $N$ is a random variable $V_N$ on $\Omega$:
the amount it pays at time $N$. It may depend only on $S_N$, as the call $(S_N-K)^+$ and
the put $(K-S_N)^+$, or on the whole path, as the lookback $(\max_{n\le N}S_n-K)^+$.

\begin{definition}[Self-financing portfolio]\label{def:bin-portfolio}
A \emph{portfolio} consists of an initial capital $X_0\in\R$ and a process
$(\Delta_n)_{0\le n\le N-1}$ adapted to $(\F_n)$: $\Delta_n$ is the number of shares held
from time $n$ to time $n+1$. Its value process is defined by
\begin{equation}\label{eq:bin-wealth}
  X_{n+1}=\Delta_nS_{n+1}+(1+r)\,(X_n-\Delta_nS_n),\qquad 0\le n\le N-1 .
\end{equation}
\end{definition}

At time $n$ the investor splits the wealth $X_n$ into $\Delta_nS_n$ in the stock and
$X_n-\Delta_nS_n$ in the money market; a negative $\Delta_n$ is a short sale and a negative
cash position a loan. The equation says that one period later the wealth is the value of
these two positions, with no money added or withdrawn: the portfolio finances itself.
Dividing \eqref{eq:bin-wealth} by $(1+r)^{n+1}$,
\begin{equation}\label{eq:bin-discounted-wealth}
  (1+r)^{-(n+1)}X_{n+1}-(1+r)^{-n}X_n=\Delta_n\Big((1+r)^{-(n+1)}S_{n+1}-(1+r)^{-n}S_n\Big).
\end{equation}
The discounted wealth is $X_0$ plus the gains of the strategy in the discounted stock: the
martingale transform of \cref{def:mg-transform}, with the predictable position
$H_{n+1}=\Delta_n$ held over the period from $n$ to $n+1$.

An \emph{arbitrage} is a portfolio with $X_0=0$, $X_N\ge0$ almost surely and
$\Prob(X_N>0)>0$: type (b) of \cref{def:arb-arbitrage}, for a portfolio rebalanced over
time (type (a) reduces to it by investing $-X_0>0$ in the money market). As every outcome
has positive probability, this means $X_N\ge0$ everywhere and $X_N>0$ somewhere; $p$ plays
no role.

\section{One period}\label{sec:bin-one-period}

In the one-period model, a portfolio is a pair of numbers $(X_0,\Delta_0)$ and its value
at time $1$ is
\begin{equation}\label{eq:bin-one-step}
  X_1=(1+r)X_0+\Delta_0\big(S_1-(1+r)S_0\big).
\end{equation}

\begin{proposition}[No arbitrage in one period]\label{prop:bin-no-arbitrage}
The one-period binomial model admits no arbitrage if and only if $d<1+r<u$.
\end{proposition}

\begin{proof}
With $X_0=0$, \eqref{eq:bin-one-step} gives $X_1(H)=\Delta_0(u-1-r)S_0$ and
$X_1(T)=\Delta_0(d-1-r)S_0$.

Suppose $d<1+r<u$, so that the first bracket is positive and the second negative. If
$\Delta_0>0$ then $X_1(T)<0$; if $\Delta_0<0$ then $X_1(H)<0$; if $\Delta_0=0$ then
$X_1=0$. No choice is an arbitrage.

Conversely, if $d\ge1+r$, buy one share with borrowed money, $\Delta_0=1$: then
$X_1(T)=(d-1-r)S_0\ge0$ and $X_1(H)=(u-1-r)S_0>0$ because $u>d$. If $u\le1+r$, sell one
share short and invest the proceeds, $\Delta_0=-1$: then $X_1(H)=(1+r-u)S_0\ge0$ and
$X_1(T)=(1+r-d)S_0>0$. Both are arbitrages.
\end{proof}

The condition says that neither the stock nor the money market dominates the other in the
sense of \cref{prop:arb-dominance}.

\begin{theorem}[Replication in one period]\label{thm:bin-one-period}
Assume $d<1+r<u$ and set
\begin{equation}\label{eq:bin-q}
  q_u=\frac{1+r-d}{u-d},\qquad q_d=\frac{u-(1+r)}{u-d}.
\end{equation}
Then $q_u,q_d\in(0,1)$ and $q_u+q_d=1$. Every claim $V_1$ is replicated by exactly one
portfolio, namely
\begin{equation}\label{eq:bin-replication}
  \Delta_0=\frac{V_1(H)-V_1(T)}{S_1(H)-S_1(T)},\qquad
  X_0=\frac{q_uV_1(H)+q_dV_1(T)}{1+r}.
\end{equation}
The number $V_0:=X_0$ is the only price of the claim at time $0$ that does not create an
arbitrage when the claim is traded together with the stock and the money market.
\end{theorem}

\begin{proof}
$q_u>0$ because $1+r>d$, $q_d>0$ because $u>1+r$, and $q_u+q_d=(u-d)/(u-d)=1$.

By \eqref{eq:bin-one-step}, $(X_0,\Delta_0)$ replicates $V_1$ if and only if
\[
  (1+r)X_0+\Delta_0(u-1-r)S_0=V_1(H),\qquad (1+r)X_0+\Delta_0(d-1-r)S_0=V_1(T).
\]
Subtracting the second equation from the first gives $\Delta_0(u-d)S_0=V_1(H)-V_1(T)$,
which determines $\Delta_0$ because $u\ne d$; this is the first formula, as
$S_1(H)-S_1(T)=(u-d)S_0$. Substituting into the first equation and using
$(u-1-r)/(u-d)=q_d$,
\[
  (1+r)X_0=V_1(H)-q_d\big(V_1(H)-V_1(T)\big)=q_uV_1(H)+q_dV_1(T).
\]
So there is at most one replicating portfolio. It does replicate: with
$(d-1-r)/(u-d)=-q_u$, the left side of the second equation equals
$q_uV_1(H)+q_dV_1(T)-q_u\big(V_1(H)-V_1(T)\big)=V_1(T)$.

Suppose the claim can be bought or sold at time $0$ for a price $c$. If $c>V_0$, sell the
claim, invest $V_0$ in the replicating portfolio and $c-V_0$ in the money market: at time
$1$ the portfolio pays $V_1$, which settles the claim, and $(1+r)(c-V_0)>0$ is left in
every state. If $c<V_0$, buy the claim, run the replicating portfolio in reverse
($-\Delta_0$ shares) and invest $V_0-c$. Either way there is an arbitrage; this is the law
of one price of \cref{prop:arb-dominance}. Conversely, the probability $\Q$ with
$\Q(H)=q_u$ and $\Q(T)=q_d$ is a pricing measure (\cref{def:arb-pricing-measure}) for the
money market and the stock, with $B=(1+r)^{-1}$: indeed $\EQ[S_1]=(q_uu+q_dd)S_0$ and
$(1+r-d)u+(u-1-r)d=(1+r)(u-d)$, so $\EQ[S_1]/(1+r)=S_0$. Since $V_0=\EQ[V_1]/(1+r)$, \cref{lem:arb-pricing-measure} shows that
adding the claim at the price $V_0$ creates no arbitrage.
\end{proof}

\subsection{The risk-neutral measure}

The weights $q_u$ and $q_d$ are positive and add up to one: they are probabilities, and
\eqref{eq:bin-replication} says that the price is a discounted expectation under them.
This is the central object of the chapter, and we define it directly for $N$ periods.

\begin{definition}[Risk-neutral measure]\label{def:bin-risk-neutral}
A probability measure $\Q$ on $\Omega$ is \emph{risk-neutral}, or an \emph{equivalent
martingale measure}, for the $N$-period binomial model if
\begin{enumerate}[label=(\roman*)]
  \item $\Q$ is equivalent to $\Prob$ (\cref{def:pr-equivalent}); since every outcome has
    positive $\Prob$-probability, this means $\Q(\omega):=\Q(\{\omega\})>0$ for every
    $\omega\in\Omega$;
  \item the discounted stock price $\big((1+r)^{-n}S_n\big)_{0\le n\le N}$ is a
    $\Q$-martingale with respect to $(\F_n)$ (\cref{def:mg-martingale}).
\end{enumerate}
\end{definition}

For $N=1$ this is a pricing measure in the sense of \cref{def:arb-pricing-measure} for the
money market and the stock. Condition (ii) reads $\EQ[S_1]=(1+r)S_0$, that is,
$\Q(H)\,uS_0+\big(1-\Q(H)\big)dS_0=(1+r)S_0$, whose only solution is $\Q(H)=q_u$; and (i)
holds if and only if $q_u\in(0,1)$, that is, $d<1+r<u$. So under the no-arbitrage
condition there is exactly one risk-neutral measure, $\Q(H)=q_u$, $\Q(T)=q_d$, and
\cref{thm:bin-one-period} reads
\begin{equation}\label{eq:bin-rn-one}
  V_0=\EQ\Big[\frac{V_1}{1+r}\Big],\qquad \Delta_0=\frac{V_1(H)-V_1(T)}{S_1(H)-S_1(T)} .
\end{equation}
The price is the expected discounted payoff, computed with $\Q$ and not with $\Prob$, and
the hedge ratio is a discrete derivative of the claim with respect to the stock. Under
$\Q$ the stock earns on average the risk-free rate, the return that an investor indifferent
to risk would accept; hence the name, which describes the measure, not the investors.

\begin{remark}[Why $p$ does not enter the price]\label{rem:bin-p}
The real-world probability $p$ appears nowhere in \eqref{eq:bin-replication}: replication
must work in every state, so only which states are possible matters. This does not make $p$
irrelevant: with investors' attitudes to risk, it shapes the stock price $S_0$, and claims
are priced \emph{relative to} $S_0$. Traders who agree on $S_0$, $u$, $d$ and $r$ but not on
$p$ must agree on every price, or one offers the other an arbitrage. What $p$ changes is
expected returns: under $\Prob$ the stock's mean return $pu+(1-p)d-1$ may exceed $r$ (a
risk premium), and a claim's in general differs from both (\exref{ex:bin-one-period}).
\end{remark}


\section{\texorpdfstring{$N$}{N} periods}\label{sec:bin-multi-period}

In $N$ periods the one-period argument is applied at every node, from the last date
backwards. The successors of a node $\omega_1\cdots\omega_n$ are $\omega_1\cdots\omega_nH$
and $\omega_1\cdots\omega_nT$.

\begin{theorem}[Replication by backward induction]\label{thm:bin-multi-period}
Assume $d<1+r<u$ and let $V_N$ be a claim. Define $V_n$ and $\Delta_n$ for
$n=N-1,\dots,0$, node by node, by
\begin{align}
  V_n(\omega_1\cdots\omega_n)&=\frac{q_uV_{n+1}(\omega_1\cdots\omega_nH)+q_dV_{n+1}(\omega_1\cdots\omega_nT)}{1+r},
  \label{eq:bin-backward}\\
  \Delta_n(\omega_1\cdots\omega_n)&=\frac{V_{n+1}(\omega_1\cdots\omega_nH)-V_{n+1}(\omega_1\cdots\omega_nT)}
  {S_{n+1}(\omega_1\cdots\omega_nH)-S_{n+1}(\omega_1\cdots\omega_nT)}.\label{eq:bin-delta}
\end{align}
\begin{enumerate}[label=(\roman*)]
  \item The portfolio with initial capital $X_0=V_0$ and hedge $(\Delta_n)$ satisfies
    $X_n=V_n$ for every $n$; in particular $X_N=V_N$ in every state. It is the only
    portfolio that replicates $V_N$.
  \item Let $\Q$ be the probability under which the tosses are independent with
    $\Q(\omega_n=H)=q_u$, that is, $\Q(\omega)=q_u^{\#H(\omega)}q_d^{\#T(\omega)}$. Then
    \begin{equation}\label{eq:bin-price}
      V_n=\EQ\big[(1+r)^{-(N-n)}V_N\,\big|\,\F_n\big],\quad 0\le n\le N;\qquad
      V_0=\EQ\big[(1+r)^{-N}V_N\big].
    \end{equation}
\end{enumerate}
\end{theorem}

\begin{proof}
(i) The formulas define $V_n$ and $\Delta_n$ as functions of $\omega_1\cdots\omega_n$, so
$(\Delta_n)$ is adapted. We show $X_n=V_n$ by induction on $n$; it holds for $n=0$ by
construction. Suppose $X_n=V_n$ at a node $\omega_1\cdots\omega_n$. Seen from this node,
the next period is a one-period model with initial price $S_n(\omega_1\cdots\omega_n)$,
the same factors $u$ and $d$, and a claim paying $V_{n+1}(\omega_1\cdots\omega_nH)$ or
$V_{n+1}(\omega_1\cdots\omega_nT)$. Formulas \eqref{eq:bin-backward} and
\eqref{eq:bin-delta} are \eqref{eq:bin-replication} for this model, so
\cref{thm:bin-one-period} gives $X_{n+1}=V_{n+1}$ at both successors.

Conversely, let $(X'_0,\Delta')$ replicate $V_N$. At each node at time $N-1$, the pair
$(X'_{N-1},\Delta'_{N-1})$ replicates $V_N$ over the last period, so the uniqueness part
of \cref{thm:bin-one-period} forces $X'_{N-1}=V_{N-1}$ and $\Delta'_{N-1}=\Delta_{N-1}$.
Repeating the argument at times $N-2,\dots,0$ gives $X'_0=V_0$ and $\Delta'=\Delta$.

(ii) Let $Y=f(\omega_1\cdots\omega_{n+1})$ be $\F_{n+1}$-measurable, and let $A$ be the
event that the first $n$ tosses are $\omega_1\cdots\omega_n$. Then $\Q(A)>0$ and
$\Q(A\cap\{\omega_{n+1}=H\})=\Q(A)\,q_u$, so the formula for conditional expectations
given a finite partition (\cref{def:pr-condexp-discrete}) gives, on $A$,
\begin{equation}\label{eq:bin-one-step-condexp}
  \EQ[Y\mid\F_n]=\frac{\EQ[Y\1_A]}{\Q(A)}=q_u\,f(\omega_1\cdots\omega_nH)+q_d\,f(\omega_1\cdots\omega_nT).
\end{equation}
(This is the independence lemma, \cref{lem:pr-independence}, since $\omega_{n+1}$ is
independent of $\F_n$ under $\Q$.) Thus \eqref{eq:bin-backward} says
$(1+r)^{-n}V_n=\EQ\big[(1+r)^{-(n+1)}V_{n+1}\mid\F_n\big]$ for every $n$, and the tower
property (\cref{prop:pr-condexp-properties}) gives
$(1+r)^{-n}V_n=\EQ\big[(1+r)^{-N}V_N\mid\F_n\big]$.
\end{proof}

As in \cref{thm:bin-one-period}, a price $c\ne V_0$ is exploited by trading the claim
against its replicating portfolio, so $V_0$ is the only candidate price, and it creates no
arbitrage (\cref{rem:bin-extended}); likewise $V_n$ is the price at time $n$.

\begin{figure}[ht]
  \centering
  \includegraphics[width=0.86\textwidth]{ch04-tree}
  \caption{Three periods with $S_0=100$, $u=1.2$, $d=0.9$ and $r=5\%$ (so $q_u=\frac12$),
    and a call with strike $100$. At each node: the stock price $S_n$, the price $V_n$ of
    the call from \eqref{eq:bin-backward}, and the number of shares $\Delta_n$ held in the
    replicating portfolio over the next period.}
  \label{fig:bin-tree}
\end{figure}

\paragraph{Recombining trees.} If $V_N=g(S_N)$, then by induction $V_n=v_n(S_n)$, with
$v_N=g$ and $v_n(s)=\big(q_uv_{n+1}(us)+q_dv_{n+1}(ds)\big)/(1+r)$, and the computation
runs over the $n+1$ values of $S_n$ rather than the $2^n$ nodes (\cref{fig:bin-tree}). Under $\Q$ the number of
heads $H_N$ is binomial with parameters $N$ and $q_u$, so \eqref{eq:bin-price} is the
\emph{Cox--Ross--Rubinstein formula}
\begin{equation}\label{eq:bin-crr}
  V_0=(1+r)^{-N}\sum_{k=0}^{N}\binom Nk q_u^k\,q_d^{N-k}\,g\big(S_0u^kd^{N-k}\big).
\end{equation}
Its limit as $N\to\infty$ is the subject of \cref{ch:06-black-scholes-formula}. A
path-dependent claim needs the whole tree (\exref{ex:bin-lookback}).

\demo{04-binomial.ipynb}{build the tree for any $u$, $d$, $r$ and $N$, read off prices and
hedge ratios at every node, check the replication along all $2^N$ paths, and see that $p$
moves expected payoffs but not prices.}

\begin{proposition}[Discounted prices are martingales]\label{prop:bin-discounted-martingale}
Assume $d<1+r<u$ and let $\Q$ be the measure of \cref{thm:bin-multi-period}.
\begin{enumerate}[label=(\roman*)]
  \item The discounted stock price $(1+r)^{-n}S_n$ is a $\Q$-martingale, so $\Q$ is a
    risk-neutral measure.
  \item Let $\Q'$ be any probability measure under which the discounted stock price is a
    martingale. For every portfolio, the discounted value $(1+r)^{-n}X_n$ is a
    $\Q'$-martingale; in particular $\E^{\Q'}\big[(1+r)^{-N}X_N\big]=X_0$.
  \item For every claim, the discounted price $(1+r)^{-n}V_n$ is a $\Q$-martingale.
\end{enumerate}
\end{proposition}

\begin{proof}
(i) $S_{n+1}$ equals $uS_n$ or $dS_n$ according to $\omega_{n+1}$, so by
\eqref{eq:bin-one-step-condexp}
\[
  \EQ[S_{n+1}\mid\F_n]=(q_uu+q_dd)\,S_n=\frac{(1+r-d)u+(u-1-r)d}{u-d}\,S_n=(1+r)\,S_n .
\]
Dividing by $(1+r)^{n+1}$ gives the martingale property, and $\Q(\omega)>0$ for every
$\omega$ because $q_u,q_d>0$.

(ii) On the finite space $\Omega$ every random variable is bounded. Write
$D_n=(1+r)^{-n}S_n$. By \eqref{eq:bin-discounted-wealth}, taking out what is known
($\Delta_n$ is $\F_n$-measurable),
\[
  \E^{\Q'}\big[(1+r)^{-(n+1)}X_{n+1}-(1+r)^{-n}X_n\,\big|\,\F_n\big]
  =\Delta_n\,\E^{\Q'}[D_{n+1}-D_n\mid\F_n]=0 .
\]
This is \cref{thm:mg-transform}: $\Delta_n$ is chosen at time $n$ and multiplies the
increment from $n$ to $n+1$. Taking expectations gives the last claim.

(iii) This is \eqref{eq:bin-price} with the tower property, or (ii) applied to the
replicating portfolio.
\end{proof}

\section{The fundamental theorem of asset pricing}\label{sec:bin-ftap}

\begin{theorem}[Fundamental theorem of asset pricing, binomial model]\label{thm:bin-ftap}
In the $N$-period binomial model the following are equivalent:
\begin{enumerate}[label=(\alph*)]
  \item there is no arbitrage;
  \item there is a risk-neutral measure;
  \item $d<1+r<u$.
\end{enumerate}
If they hold, the risk-neutral measure is unique, it is the measure $\Q$ of
\cref{thm:bin-multi-period}, and every claim can be replicated: the model is
\emph{complete}.
\end{theorem}

\begin{proof}
(b)$\Rightarrow$(a). Let $\Q'$ be risk-neutral and $X$ a portfolio with $X_0=0$ and
$X_N\ge0$. By \cref{prop:bin-discounted-martingale}(ii),
$\E^{\Q'}\big[(1+r)^{-N}X_N\big]=0$. A nonnegative random variable with mean zero
vanishes at every outcome of positive probability, and $\Q'(\omega)>0$ for every
$\omega$; so $X_N=0$ everywhere and $X$ is not an arbitrage. (This is the argument of
\cref{lem:arb-pricing-measure}, with the martingale property of discounted wealth in place
of the linearity of static prices.)

(a)$\Rightarrow$(c). Suppose $d\ge1+r$. Borrow $S_0$, buy one share at time $0$, sell it
at time $1$ and keep the balance in the money market until time $N$: $X_0=0$,
$\Delta_0=1$ and $\Delta_n=0$ for $n\ge1$. Then $X_1=S_1-(1+r)S_0$ is nonnegative, and
positive if $\omega_1=H$, as in the proof of \cref{prop:bin-no-arbitrage}; and
$X_N=(1+r)^{N-1}X_1$. Since $\Prob(\omega_1=H)=p>0$, this is an arbitrage. If $u\le1+r$,
take $\Delta_0=-1$ instead.

(c)$\Rightarrow$(b). This is \cref{prop:bin-discounted-martingale}(i).

Uniqueness and completeness. Under (c) every claim can be replicated, by
\cref{thm:bin-multi-period}. Let $\Q'$ be a risk-neutral measure and fix an outcome
$\varpi$. The claim $\1_{\{\varpi\}}$, which pays one dollar if $\omega=\varpi$ and nothing
otherwise, is replicated by a portfolio $X$ with
$X_0=\EQ\big[(1+r)^{-N}\1_{\{\varpi\}}\big]=(1+r)^{-N}\Q(\varpi)$, by \eqref{eq:bin-price}.
By \cref{prop:bin-discounted-martingale}(ii) applied to $\Q'$, also
$X_0=\E^{\Q'}\big[(1+r)^{-N}X_N\big]=(1+r)^{-N}\Q'(\varpi)$. Hence $\Q'(\varpi)=\Q(\varpi)$
for every $\varpi$.
\end{proof}

The claims $\1_{\{\varpi\}}$ are the \emph{Arrow--Debreu securities} of the model, and
their prices $(1+r)^{-N}\Q(\varpi)$ are the state prices of \cref{ch:03-arbitrage}
(\exref{ex:bin-state-prices}).

\begin{remark}[Trading the claim creates no arbitrage]\label{rem:bin-extended}
If the claim is traded at $V_n$, its discounted price is a $\Q$-martingale
(\cref{prop:bin-discounted-martingale}(iii)), and the proof of (b)$\Rightarrow$(a) applies
verbatim to portfolios in the three assets: $V_n$ is a consistent price.
\end{remark}

\begin{remark}[The general theorem]\label{rem:bin-ftap-general}
The equivalence (a)$\Leftrightarrow$(b) holds in every market with finitely many trading
dates and finitely many assets: \textcite{harrison-pliska} proved the multi-period version
on a finite probability space, and \textcite{dalang-morton-willinger} extended it to general
probability spaces; see \textcite{follmer-schied}, Section~1.2 for one period and
Section~5.2 for several. On a finite $\Omega$, (b)$\Rightarrow$(a) has the proof above; in
general the discounted wealth of an unbounded strategy need not be an integrable
martingale, and one first shows that it is when its terminal value is nonnegative
(\textcite{follmer-schied}, Chapter~5). The converse needs a separating hyperplane
argument, which the binomial model replaces by the explicit formula \eqref{eq:bin-q}. Completeness is also general: an
arbitrage-free market of this kind is complete if and only if its risk-neutral measure is
unique (\textcite{follmer-schied}, Section~5.4); the Arrow--Debreu argument above is one
half of that statement, and \cref{sec:bin-trinomial} shows the other half at work.
\end{remark}

\section{American options}\label{sec:bin-american}

An American option can be exercised at any date up to its maturity $N$; exercise at time
$n$ pays $G_n$, where $(G_n)$ is adapted: $G_n=(K-S_n)^+$ for the American put and
$G_n=(S_n-K)^+$ for the American call. The holder decides when to exercise with the
information available, so an exercise strategy is a stopping time
(\cref{def:mg-stopping-time}): a random time $\tau$ with values in $\{0,\dots,N\}$ such
that $\{\tau=n\}\in\F_n$ for every $n$. We write $\mathcal T_n$ for
the set of stopping times with $n\le\tau\le N$.

The price is computed backwards, comparing at every node exercise with continuation:
\begin{equation}\label{eq:bin-american}
  U_N=G_N,\qquad U_n=\max\Big(G_n,\ \frac{1}{1+r}\EQ[U_{n+1}\mid\F_n]\Big),\quad n=N-1,\dots,0 .
\end{equation}
The second term is the one-period price of a claim paying $U_{n+1}$ one period later: the
value of keeping the option alive. \eqref{eq:bin-american} is a dynamic programming
equation, and $(1+r)^{-n}U_n$ is the \emph{Snell envelope} of $(1+r)^{-n}G_n$.

The proof rests on optional stopping, which we need conditionally on $\F_n$; for $n=0$ the
next lemma is \cref{thm:mg-optional-stopping}(a).

\begin{lemma}[Optional stopping on a finite horizon]\label{lem:bin-optional-stopping}
On any probability space with a filtration $(\F_k)_{0\le k\le N}$, let $(Y_k)_{0\le k\le N}$
be a supermartingale and $\tau$ a stopping time with $n\le\tau\le N$. Then
$\E[Y_\tau\mid\F_n]\le Y_n$, with equality if $Y$ is a martingale.
\end{lemma}

\begin{proof}
$Y_\tau=\sum_k\1_{\{\tau=k\}}Y_k$ is integrable. Since $n\le\tau\le N$, the following sum
telescopes:
\[
  Y_\tau-Y_n=\sum_{k=n+1}^{N}\1_{\{\tau\ge k\}}\,(Y_k-Y_{k-1}).
\]
The event $\{\tau\ge k\}=\Omega\setminus\{\tau\le k-1\}$ belongs to $\F_{k-1}$. For
$k\ge n+1$, the tower property, taking out what is known and positivity give
\[
  \E\big[\1_{\{\tau\ge k\}}(Y_k-Y_{k-1})\,\big|\,\F_n\big]
  =\E\big[\1_{\{\tau\ge k\}}\,\E[Y_k-Y_{k-1}\mid\F_{k-1}]\,\big|\,\F_n\big]\le0,
\]
with equality for a martingale. Summing over $k$ proves the lemma.
\end{proof}

\begin{theorem}[American options]\label{thm:bin-american}
Assume $d<1+r<u$, let $G$ be an adapted payoff process, define $U$ by
\eqref{eq:bin-american}, and let $\tau^*_n=\min\{k\ge n:U_k=G_k\}$.
\begin{enumerate}[label=(\roman*)]
  \item $U_n\ge G_n$ for every $n$, and $(1+r)^{-n}U_n$ is a $\Q$-supermartingale.
  \item $\tau^*_n\in\mathcal T_n$, and for every $\tau\in\mathcal T_n$,
    \[
      (1+r)^n\,\EQ\big[(1+r)^{-\tau}G_\tau\,\big|\,\F_n\big]\le U_n
      =(1+r)^n\,\EQ\big[(1+r)^{-\tau^*_n}G_{\tau^*_n}\,\big|\,\F_n\big].
    \]
    In particular $U_0=\max_{\tau\in\mathcal T_0}\EQ\big[(1+r)^{-\tau}G_\tau\big]$, and the
    maximum is attained at $\tau^*:=\tau^*_0$.
  \item There is a portfolio with initial capital $U_0$ whose value satisfies
    $X_n\ge U_n\ge G_n$ for every $n$; and every portfolio with $X_n\ge G_n$ for all $n$
    has $X_0\ge U_0$.
  \item (American call) Let $r\ge0$, $K\ge0$ and $G_n=(S_n-K)^+$, and let $V_n$ be the
    price \eqref{eq:bin-price} of the European call with payoff $(S_N-K)^+$. Then $U_n=V_n$
    for every $n$, and $\tau=N$ is an optimal exercise time: nothing is gained by exercising
    early. If $r>0$ and $K>0$, exercise at a time $n<N$ in a state where $S_n>K$ is
    strictly suboptimal: $U_n>G_n$ there.
\end{enumerate}
\end{theorem}

\begin{proof}
Write $Y_k=(1+r)^{-k}U_k$ and $c_k=\EQ[U_{k+1}\mid\F_k]/(1+r)$ for the continuation value.

(i) Both terms of the maximum in \eqref{eq:bin-american} are $\F_n$-measurable, so $U$ is
adapted. By definition $U_n\ge G_n$, and $U_n\ge c_n$ is, after division by $(1+r)^n$,
the supermartingale inequality $Y_n\ge\EQ[Y_{n+1}\mid\F_n]$.

(ii) Since $U_N=G_N$, we have $n\le\tau^*_n\le N$, and for $n\le k\le N$,
$\{\tau^*_n=k\}=\{U_n>G_n,\dots,U_{k-1}>G_{k-1},U_k=G_k\}\in\F_k$; so
$\tau^*_n\in\mathcal T_n$. Let $\tau\in\mathcal T_n$. By \cref{lem:bin-optional-stopping}
for the supermartingale $Y$, and by (i),
\[
  Y_n\ge\EQ[Y_\tau\mid\F_n]\ge\EQ\big[(1+r)^{-\tau}G_\tau\,\big|\,\F_n\big],
\]
which is the inequality. For $\tau^*_n$, consider the stopped process
$M_k=Y_{k\wedge\tau^*_n}$, $n\le k\le N$. Its increments are
$M_{k+1}-M_k=\1_{\{\tau^*_n>k\}}(Y_{k+1}-Y_k)$, and $\{\tau^*_n>k\}\in\F_k$. On this event
$U_k>G_k$, so $U_k=c_k$, that is, $\EQ[Y_{k+1}\mid\F_k]=Y_k$. Taking out what is known,
\[
  \EQ[M_{k+1}-M_k\mid\F_k]=\1_{\{\tau^*_n>k\}}\big(\EQ[Y_{k+1}\mid\F_k]-Y_k\big)=0,
\]
so $M$ is a $\Q$-martingale on $\{n,\dots,N\}$ and, by the tower property,
$Y_n=M_n=\EQ[M_N\mid\F_n]=\EQ[Y_{\tau^*_n}\mid\F_n]$. Finally
$Y_{\tau^*_n}=(1+r)^{-\tau^*_n}G_{\tau^*_n}$, because $U=G$ at time $\tau^*_n$.

(iii) Define $\Delta_n$ by \eqref{eq:bin-delta} with $U$ in place of $V$, and start from
$X_0=U_0$. At each node, \cref{thm:bin-one-period} says that the capital $c_n$ with
$\Delta_n$ shares replicates $U_{n+1}$ one period later:
$(1+r)c_n+\Delta_n\big(S_{n+1}-(1+r)S_n\big)=U_{n+1}$. Subtracting this from
\eqref{eq:bin-wealth} written as $X_{n+1}=(1+r)X_n+\Delta_n(S_{n+1}-(1+r)S_n)$,
\[
  X_{n+1}=U_{n+1}+(1+r)(X_n-c_n).
\]
If $X_n\ge U_n$, then $X_n\ge c_n$ and so $X_{n+1}\ge U_{n+1}$; by induction $X_n\ge U_n$
for every $n$. Conversely, let $X$ be a portfolio with $X_n\ge G_n$ for all $n$. Its
discounted value is a $\Q$-martingale (\cref{prop:bin-discounted-martingale}(ii)), so
\cref{lem:bin-optional-stopping} with $\tau=\tau^*$ gives
$X_0=\EQ\big[(1+r)^{-\tau^*}X_{\tau^*}\big]\ge\EQ\big[(1+r)^{-\tau^*}G_{\tau^*}\big]=U_0$.

(iv) By the conditional Jensen inequality for the convex function $x\mapsto(x-K)^+$
(\cref{prop:pr-condexp-properties}) and the martingale property
$\EQ[S_N\mid\F_n]=(1+r)^{N-n}S_n$,
\begin{equation}\label{eq:bin-call-bound}
\begin{split}
  V_n=\frac{\EQ[(S_N-K)^+\mid\F_n]}{(1+r)^{N-n}}
  &\ge\frac{\big((1+r)^{N-n}S_n-K\big)^+}{(1+r)^{N-n}}\\
  &=\big(S_n-(1+r)^{-(N-n)}K\big)^+\ge(S_n-K)^+=G_n,
\end{split}
\end{equation}
where the last step uses $(1+r)^{-(N-n)}K\le K$, valid because $r\ge0$ and $K\ge0$. Now
$U_N=G_N=V_N$, and if $U_{n+1}=V_{n+1}$, then \eqref{eq:bin-american} and
\eqref{eq:bin-backward} give $U_n=\max(G_n,V_n)=V_n$. By backward induction $U_n=V_n$ for
all $n$, and $\tau=N$ attains $U_0=V_0=\EQ\big[(1+r)^{-N}G_N\big]$. If $r>0$, $K>0$,
$n<N$ and $S_n>K$, \eqref{eq:bin-call-bound} gives
$U_n=V_n\ge S_n-(1+r)^{-(N-n)}K>S_n-K=G_n$.
\end{proof}

\paragraph{The arbitrage-free price.} Suppose the option can be bought or sold at time $0$
for $c$, and trading in the stock and the money market is allowed. The buyer chooses an
exercise time $\tau\in\mathcal T_0$ and keeps $G_\tau$ in the money market; the seller pays
$G_\tau$ at whatever $\tau$ the buyer chooses. Following \textcite{follmer-schied},
Chapter~6, call $c$ \emph{arbitrage-free} if (a) for no $\tau$ can the buyer, borrowing $c$,
exercising at $\tau$ and trading, end at time $N$ with a nonnegative and not identically
zero amount; and (b) for some $\tau$ the seller, receiving $c$, paying $G_\tau$ and trading,
cannot. Then $U_0$ is the only arbitrage-free price. If $c>U_0$, the seller invests $U_0$ in
the portfolio of (iii) and $c-U_0$ in the money market; whatever $\tau$, $X_\tau\ge G_\tau$
covers the payment, and (b) fails. If $c<U_0$, the buyer exercises at $\tau^*$ and sells
short the portfolio replicating $(1+r)^{N-\tau^*}G_{\tau^*}$ paid at $N$, worth $U_0$ by
\eqref{eq:bin-price} and (ii), and keeps $U_0-c>0$; (a) fails. If $c=U_0$, a buyer using
$\tau$ and a portfolio started from $0$ ends with an amount whose discounted $\Q$-mean is
$\EQ[(1+r)^{-\tau}G_\tau]-U_0\le0$ by (ii) and \cref{prop:bin-discounted-martingale}(ii);
a seller facing $\tau^*$ ends with discounted $\Q$-mean $U_0-\EQ[(1+r)^{-\tau^*}G_{\tau^*}]=0$.
As $\Q(\omega)>0$ for every $\omega$, a nonnegative amount with nonpositive mean is zero:
neither has an arbitrage. The time $\tau^*$ is the smallest optimal exercise time; there may
be others (\exref{ex:bin-snell}).

In (iv), the bound \eqref{eq:bin-call-bound} is the lower bound of \cref{prop:arb-bounds}
at time $n$, which holds in every arbitrage-free market, as \textcite{merton1973} observed:
a call on a stock that pays no dividend is worth more alive than exercised. (The time $\tau^*$ may still stop at nodes where the call is
worthless, $U_n=G_n=0$, which changes nothing.) For the put the argument fails: the corresponding bound $V_n\ge\big((1+r)^{-(N-n)}K-S_n\big)^+$
lies below the exercise value $K-S_n$ when $r>0$, and deep in the money early exercise is
strictly better (\cref{fig:bin-american-put}). With $r=0$ the American and European puts
coincide (\exref{ex:bin-put}).

\begin{figure}[ht]
  \centering
  \includegraphics[width=0.9\textwidth]{ch04-american-put}
  \caption{American and European puts with $K=100$ on a tree with $N=40$, $u=1.04$,
    $d=1/u$ and $r=0.3\%$ per period. Left: prices at time $0$ as functions of $S_0$; the
    American price meets the payoff where immediate exercise is optimal. Right: the nodes
    of the tree started at $S_0=K$, marked where exercise is optimal ($G_n>0$ and $G_n$ at
    least the continuation value).}
  \label{fig:bin-american-put}
\end{figure}

\demo{04-binomial.ipynb}{compare American and European puts and watch the early-exercise
region as $r$, $u$, $N$ and the strike move.}

\section{The trinomial model: a first incomplete market}\label{sec:bin-trinomial}

In the one-period \emph{trinomial model}, $S_1$ takes the values $uS_0$, $mS_0$, $dS_0$,
with $0<d<m<u$ and positive $\Prob$-probability for each, and the money market pays
$R:=1+r$ per dollar. A risk-neutral measure $(q_u,q_m,q_d)$ must satisfy two linear
equations in three unknowns. A portfolio is worth $RX_0+\Delta_0(S_1-RS_0)$, affine in
$S_1$, so it replicates only claims whose three values lie on a line; a call with strike
$dS_0<K<uS_0$ is not one of them. We call $c$ an
\emph{arbitrage-free price} for a claim $h(S_1)$ if the market in which the money market,
the stock and the claim (at price $c$ at time $0$) are traded admits no arbitrage.

\begin{proposition}[The one-period trinomial model]\label{prop:bin-trinomial}
Let $0<d<m<u$ and $R=1+r$.
\begin{enumerate}[label=(\roman*)]
  \item The model has no arbitrage if and only if $d<R<u$.
  \item In that case the risk-neutral measures are the measures $\Q_\theta$,
    $0<\theta<\theta_{\max}$, given by
    \[
      q_m=\theta,\quad q_u=\frac{R-d-\theta(m-d)}{u-d},\quad
      q_d=\frac{u-R-\theta(u-m)}{u-d},\quad
      \theta_{\max}=\frac{R-d}{m-d}\wedge\frac{u-R}{u-m}.
    \]
  \item Assume $d<R<u$, let $\theta_{\max}$ be as in (ii), let $dS_0<K<uS_0$ and
    $h(x)=(x-K)^+$, and write $h_x=h(xS_0)$ for $x\in\{d,m,u\}$. A price $c$ for the
    call is arbitrage-free if and only if
    $C_{\rm low}<c<C_{\rm high}$, where
    \[
      C_{\rm high}=\frac1R\Big(\frac{R-d}{u-d}\,h_u+\frac{u-R}{u-d}\,h_d\Big),\quad
      C_{\rm low}=\begin{cases}
        \dfrac1R\Big(\dfrac{R-m}{u-m}\,h_u+\dfrac{u-R}{u-m}\,h_m\Big) & \text{if } m\le R,\\[2ex]
        \dfrac1R\Big(\dfrac{R-d}{m-d}\,h_m+\dfrac{m-R}{m-d}\,h_d\Big) & \text{if } m\ge R.
      \end{cases}
    \]
    The interval $(C_{\rm low},C_{\rm high})$ is the set
    $\{\E^{\Q_\theta}[h(S_1)]/R:0<\theta<\theta_{\max}\}$.
\end{enumerate}
\end{proposition}

\begin{proof}
(i) If $d<R<u$, then $\theta_{\max}>0$ and the measures of (ii) exist. Each of them is a
pricing measure (\cref{def:arb-pricing-measure}) for the money market and the stock, with
$B=1/R$, because $\EQ[S_1]=RS_0$; so there is no arbitrage, by
\cref{lem:arb-pricing-measure}. If $d\ge R$, buying a share with borrowed money
gives $X_1=S_1-RS_0\ge0$, positive in the states $m$ and $u$; if $u\le R$, selling one
short gives $X_1=RS_0-S_1\ge0$, positive in the states $d$ and $m$.

(ii) The conditions are $q_u+q_m+q_d=1$ and $uq_u+mq_m+dq_d=R$. For $q_m=\theta$ fixed,
this is a linear system for $(q_u,q_d)$ with nonzero determinant $d-u$, and substitution
checks that the displayed values solve it. They are positive exactly when
$0<\theta<\theta_{\max}$.

(iii) The price under $\Q_\theta$, $\pi(\theta)=(q_uh_u+\theta h_m+q_dh_d)/R$, is affine
in $\theta$:
\[
  R\,\pi(\theta)=R\,\pi(0)-\theta\,(\ell_m-h_m),\qquad
  \ell_m=\frac{(m-d)\,h_u+(u-m)\,h_d}{u-d},
\]
where $\ell_m$ is the value at $mS_0$ of the line $\ell$ through $(dS_0,h_d)$ and
$(uS_0,h_u)$. The payoff is convex, so $h_m\le\ell_m$, and here strictly: $h_d=0<h_u$, so
if $mS_0\le K$ then $h_m=0<\ell_m$, and if $mS_0>K$ a short computation gives
$\ell_m-h_m=(u-m)(K-dS_0)/(u-d)>0$. Hence $\pi$ is strictly decreasing and
$\{\pi(\theta):0<\theta<\theta_{\max}\}=\big(\pi(\theta_{\max}),\pi(0)\big)$. The
measure $\Q_0$ has $q_m=0$ and is the binomial measure on $\{d,u\}$, so
$\pi(0)=C_{\rm high}$. One checks that $(u-R)/(u-m)\le(R-d)/(m-d)$ if and only if
$m\le R$. If $m\le R$, then $\theta_{\max}=(u-R)/(u-m)$, where $q_d=0$ and
$q_u=(R-m)/(u-m)$, and $\pi(\theta_{\max})$ is the first formula for $C_{\rm low}$; if
$m\ge R$, then $\theta_{\max}=(R-d)/(m-d)$, where $q_u=0$, and we get the second.

Arbitrage-free prices. If $C_{\rm low}<c<C_{\rm high}$, then $c=\pi(\theta)$ for some
$\theta\in(0,\theta_{\max})$: the call is priced at $\E^{\Q_\theta}[h(S_1)]/R$, and
\cref{lem:arb-pricing-measure} shows that adding it to the market creates no arbitrage. For
the converse, note that an affine payoff $\alpha+\beta S_1$ costs $\alpha/R+\beta S_0$, which
is $\E^{\Q^*}[\alpha+\beta S_1]/R$ for every probability $\Q^*$ on the three states with
$\E^{\Q^*}[S_1]=RS_0$, including $\Q_0$ and $\Q_{\theta_{\max}}$.

Suppose $c\ge C_{\rm high}$. The payoff $\ell(S_1)$ costs
$\E^{\Q_0}[\ell(S_1)]/R=\E^{\Q_0}[h(S_1)]/R=C_{\rm high}$, because $\ell=h$ at the two
states that carry $\Q_0$. Sell the call at $c$, buy $\ell(S_1)$ and invest
$c-C_{\rm high}$: the position costs nothing, and at time $1$ it is worth
$\ell(S_1)-h(S_1)+R(c-C_{\rm high})$, which is nonnegative in every state and at least
$\ell_m-h_m>0$ in the state $m$. This is an arbitrage.

Suppose $c\le C_{\rm low}$ and $m\le R$. Let $\ell'$ be the affine function through
$(mS_0,h_m)$ and $(uS_0,h_u)$; as above, $\ell'(S_1)$ costs
$\E^{\Q_{\theta_{\max}}}[h(S_1)]/R=C_{\rm low}$. Write $mS_0=\lambda dS_0+(1-\lambda)uS_0$
with $\lambda\in(0,1)$; convexity of $h$ and affinity of $\ell'$ give
$\lambda\ell'(dS_0)+(1-\lambda)h_u=h_m\le\lambda h_d+(1-\lambda)h_u$, so
$\ell'(dS_0)\le h_d$, and the inequality is strict, since otherwise the three points
$(xS_0,h_x)$ would lie on $\ell'$ and $\ell_m=h_m$. Buy the call at $c$, sell
$\ell'(S_1)$ and invest $C_{\rm low}-c$: at time $1$ the position is worth
$h(S_1)-\ell'(S_1)+R(C_{\rm low}-c)\ge0$, and positive in the state $d$. The case $m\ge R$
is symmetric, with the line through $(dS_0,h_d)$ and $(mS_0,h_m)$ and the state $u$.
\end{proof}

No arbitrage now pins down the price only up to an interval. $C_{\rm high}$ is the
\emph{superreplication price}: $\ell(S_1)$ dominates the call, and every dominating portfolio
costs $\E^{\Q_\theta}[X_1]/R\ge\pi(\theta)$ for all $\theta$; $C_{\rm low}$ is the
subreplication price. Choosing a price inside needs more information, such as the price of
another option; \exref{ex:bin-trinomial} shows this, and that a claim is replicable exactly
when its price is the same under every risk-neutral measure (cf.\ \cref{rem:bin-ftap-general}).

\begin{figure}[ht]
  \centering
  \includegraphics[width=0.9\textwidth]{ch04-trinomial}
  \caption{Arbitrage-free prices of a call with $K=S_0$ in the one-period trinomial model
    with $u=1.2$, $d=0.8$, $r=5\%$. Left: the interval of \cref{prop:bin-trinomial} as $m$
    moves from $d$ to $u$. For $K/S_0\le m\le1+r$ the lower end is the static bound
    $S_0-K/(1+r)$ of \cref{prop:arb-bounds}: the payoff is affine on the two states that
    carry $\Q_{\theta_{\max}}$. Right: for $m=1$, the price is affine in $\theta=q_m$; the
    excluded end points give zero weight to one state.}
  \label{fig:bin-trinomial}
\end{figure}

\demo{04-binomial.ipynb}{move the middle factor $m$, the strike and the rate, and watch
the interval of arbitrage-free call prices.}

\begin{remark}[Trinomial trees as numerical schemes]\label{rem:bin-trinomial-trees}
Trinomial trees used for computation (\cref{ch:16-numerics}) choose one measure by matching
the mean $\nu\,\delta t$ and second moment $s^2:=\sigma^2\delta t+\nu^2\delta t^2$ of a
log-price step $\pm\delta x$ or $0$ (the Black--Scholes values of
\cref{ch:06-black-scholes-formula}): $p_{u,d}=\frac12\big(s^2/\delta x^2\pm\nu\,\delta t/\delta x\big)$,
$p_m=1-s^2/\delta x^2$. $\delta x=\sigma\sqrt{\delta t}$ gives $p_m=-\nu^2\delta t/\sigma^2<0$
if $\nu\ne0$; $\delta x=\sigma\sqrt{3\,\delta t}$ gives $p_m=\frac23-\nu^2\delta t/(3\sigma^2)>0$
for $\delta t<2\sigma^2/\nu^2$ and $p_{u,d}=(1+a^2\pm\sqrt3\,a)/6>0$, $a=\nu\sqrt{\delta t}/\sigma$.
Such a tree is only approximately risk-neutral: $\E[e^{\delta X}]\ne e^{r\delta t}$, with a
per-step error of order $\delta t^2$ in general and $\delta t^3$ for $\delta x=\sigma\sqrt{3\,\delta t}$.
\end{remark}

\furtherreading
\textcite{shreve1}, Chapters~1--4, covers this material at the same level; the model is due
to \textcite{crr1979}. \textcite{follmer-schied} treat discrete-time markets with full
rigour: Chapters~1 and~5 for arbitrage, completeness and incomplete markets, Chapter~6 for
American options. \textcite{hull} explains how trees are used in practice.
